\section{Ordered multi-series EPS report records}
\label{sec:ordered-eps-records}

\code{OrderedEPS/1} specifies the report taught in
Sections~\ref{sec:ordered-eps-method}--\ref{sec:ordered-eps-report}.
Its closed Draft 2020-12 JSON schemas are retained as
\code{ordered-eps/record-schemas.json} in the continuation evidence directory.
Every field below is required; unknown fields reject. Braces denote objects,
brackets arrays, and alternatives distinct shapes. No migration from
\code{FinancingComponents/2} is inferred.

\begin{lstlisting}
Request {
  wire: "OrderedEPS/1", operation: "report",
  context, context_ref, inventory, scope, basic, series
}
context {
  entity, book, head, view, population,
  period: {start, end_exclusive}, ordinary_class, share_basis,
  currency, currency_scale, rule_version
}
support {path, sha256, anchor}
inventory {status: complete | incomplete, ids: [id], support}
scope {
  nonparticipating_rights, independent_additive_effects,
  fixed_effects, support
}
basic (ready) {
  status: ready, context_ref, support, continuing_numerator,
  total_numerator, denominator,
  attribution_complete, chronology_complete
}
basic (unresolved) {
  status: unresolved, context_ref, support, reason
}
\end{lstlisting}

Names and reasons are nonempty strings, findings are Booleans, dates are
valid ISO calendar dates, and intervals include their start but exclude their
end. A reporting period has positive length. Digests are 64 lower-case
hexadecimal characters. The context reference is SHA256 of UTF-8 JSON with
keys sorted, compact separators, ASCII escaping and no trailing newline.
Every basic or series record must carry that reference. A support reference
identifies exact local bytes and a nonempty anchor. Source resolution checks
the bytes; a caller must separately interpret the anchor and the finding.

Quantities are reduced rational strings: \code{0}, a signed nonzero integer,
or a reduced numerator/positive-denominator fraction. Leading zeros, plus
signs, negative zero, decimals, exponents and a redundant denominator of one
reject. The basic denominator is positive. The context binds units and scale
before arithmetic; the report performs no currency or share-basis conversion.
Both completeness findings in a ready basic record must be true to return
basic EPS. The inventory and series IDs must agree bijectively, with no
duplicates. A complete empty inventory is permitted with support.

\begin{lstlisting}
series (candidate) {
  id, resolution: candidate, context_ref, support,
  continuing_increment, total_increment, share_increment,
  earnings_complete, time_weighting_complete,
  potential_windows: [{start, end_exclusive}],
  actual_windows: [{start, end_exclusive}],
  end_kind: period_end | lapsed | converted, actual_in_basic
}
series (no effect) {
  id, resolution: no_effect, context_ref, support,
  continuing_increment, total_increment, share_increment,
  reason
}
series (unresolved or unsupported) {
  id, resolution: unresolved | unsupported,
  context_ref, support, reason
}
\end{lstlisting}

A candidate has a positive, already weighted share increment; earnings
increments may be signed. Both completeness findings must be true for diluted
calculation. It has at least one potential window. All windows have positive
length, lie within the reporting period and do not overlap within that series.
For a conversion, actual windows are nonempty, follow the potential windows,
and \code{actual\_in\_basic} is true. For lapse or period-end treatment, actual
windows are empty and that finding is false. A period-end series' final
potential window ends at the period boundary. A no-effect record has all three
increments zero. Reasons retain missing facts or the further method required;
they never stand for a silently omitted series.

\begin{lstlisting}
Result {
  wire: "OrderedEPS/1", request_sha256, context_ref,
  status, basic: EPSBundle | null, diluted: EPSBundle | null,
  trace: [Comparison], issues: [{code, path, series_id}],
  mutations: [], acceptance: "conditional_documentary_only"
}
EPSBundle {
  continuing_numerator, total_numerator,
  discontinued_numerator, denominator,
  continuing_eps, total_eps, discontinued_eps
}
Comparison {
  id, resolution: candidate | no_effect,
  continuing_increment, total_increment, share_increment,
  incremental_ratio,
  before: {continuing_numerator, denominator},
  after_if_included: {continuing_numerator, denominator},
  included, reason: strictly_dilutive | neutral |
                    antidilutive | no_effect
}
\end{lstlisting}

Result quantities remain reduced rational strings. The request digest uses
the same serialization over the complete request. The context reference is
null only when no valid reference can be read. An issue identifies its code,
JSON path and series ID or null. A no-effect comparison has a null ratio and
unchanged amounts. Candidates appear in exact increasing-ratio order, IDs
breaking ties; no-effect items follow in ID order. The schema permits null
comparison states; the completed reference report supplies them.

\code{conditional\_complete} returns both bundles, the complete trace and no
issues. \code{unresolved} records missing findings or source bytes, incomplete
inventory or unresolved items. \code{unsupported} records unsupported scope
or series; \code{blocked} records both obstacles. These three statuses retain
independently ready basic EPS, with null diluted EPS and an empty trace.
\code{rejected} takes precedence for structural, binding, rational, sign,
inventory, date or source-hash violations and returns neither bundle. Issues
sort by path, code and series ID; mutations are always empty. Application
methods remain responsible for live head freshness, source truth, accounting
admission and future publication.
